\documentclass[UTF8, fontset=macnew, a4paper, 12pt]{ctexart}

\usepackage{geometry}
\geometry{top=2.5cm, bottom=2.5cm, left=2.5cm, right=2.5cm}
\usepackage{amsmath, amssymb}
\usepackage{booktabs}
\usepackage{multirow}
\usepackage{tabularx}
\usepackage{graphicx}
\usepackage{caption}
\usepackage{xcolor}
\usepackage{float}
\usepackage{enumitem}
\usepackage{fancyhdr}
\usepackage{setspace}
\usepackage{pifont}
\usepackage{natbib}
% 拉丁作者名之後採半形括號（中文作者名在書目的中文部分仍用全形）；
% 第五個引數留空表示作者與年份之間不加逗號。
\bibpunct{(}{)}{; }{a}{}{,}
\usepackage[hidelinks]{hyperref}

\graphicspath{{../output/figures/}}
\captionsetup{font=small, labelfont=bf, labelsep=quad}
\setlist{itemsep=2pt, parsep=0pt, topsep=4pt}
\onehalfspacing

\renewcommand{\figurename}{圖}
\renewcommand{\tablename}{表}

\ctexset{
  section/format = \raggedright\zihao{4}\bfseries,
  section/name   = {},
  section/number = \bignum{\value{section}},
  section/aftername = 、,
  subsection/format = \raggedright\zihao{5}\bfseries,
  subsection/number = \chinese{subsection},
  subsection/aftername = 、,
  subsubsection/format = \raggedright\normalsize\bfseries,
}
\newcommand{\bignum}[1]{\ifcase#1\or 壹\or 貳\or 參\or 肆\or 伍\or
  陸\or 柒\or 捌\or 玖\or 拾\fi}
\renewcommand{\thesection}{\bignum{\value{section}}}
\renewcommand{\thesubsection}{\chinese{subsection}}
\makeatletter
\renewcommand{\p@subsection}{\thesection 節之}
\makeatother
\setcounter{secnumdepth}{2}

\pagestyle{fancy}
\fancyhf{}
\fancyhead[L]{}
\fancyhead[R]{\small 深入研究（一）}
\fancyfoot[C]{\small \thepage}
\renewcommand{\headrulewidth}{0.4pt}

\newcommand{\Ex}{\mathbb{E}}
\newtheorem{prop}{Proposition}
\newtheorem{cor}{Corollary}
\newtheorem{lem}{Lemma}
\newenvironment{proof}{\par\noindent\textit{證明}.\ }{\hfill$\square$\par\vspace{4pt}}

\begin{document}

\begin{center}
{\zihao{3}\bfseries 小人口 Lee-Carter 偏誤校正的深入研究（一）}\\[4pt]
{\zihao{5} ——對數轉換的二階解析量與漸近框架}\\[6pt]
{\zihao{5} Second-Order Analytic Quantities of the Log Transform and the
Asymptotic Framework}\\[12pt]
{\zihao{5} （林子立）\quad 指導教授：余清祥\,教授}\\[4pt]
{\zihao{5} 國立政治大學統計學系}\\[4pt]
{\zihao{5} 中華民國 115 年 10 月}
\end{center}

\vspace{6pt}
\begin{center}
\begin{minipage}{0.92\textwidth}
\small
\textbf{本篇已由進度報告（1005 V2）取代}\par\vspace{3pt}
本篇的全部內容已併入 \texttt{進度報告\_1005\_V2} 第肆節，並在該處作了兩項
更正，故以該版為準。其一，本篇第參節之四的變異數膨脹因子
$(1-c_x)^2$ 所對應的是\textbf{一步}校正，而本研究各節所用的估計量為
\textbf{迭代至收斂}者，後者的代價為 $1/A(\mu)^2$，$A=1+\mu b'(\mu)$，
其在零格主導的年齡發散而非趨於 $4$；實測亦支持後者
（$N=5\times10^4$ 的 5 至 9 歲組實測 $31.5$，$1/A^2$ 預測 $32.7$）。
其二，本篇第伍節的「待證命題一、二」未經證明，V2 已將該節改寫為一段說明，
指出本研究的結論建立在有限樣本恆等式之上，不依賴任何漸近敘述。
V2 另新增相對效率的 Corollary。
\par\vspace{6pt}
\textbf{摘要}\par\vspace{4pt}
本篇承接 1005 進度報告的一階結果。該報告指出標準 Lee-Carter 的
$\hat\alpha_x$ 偏誤有閉式 $b(\mu;c_0)=\Ex[\log\max(D,c_0)]-\log\mu$，
並據以設計一個不需參考母體亦不需調參的修正；但該報告只處理估計方程期望的
偏離，未處理其變異數，故留下兩項未解釋的實證現象：標準 Lee-Carter 的
$\alpha$ 帶在幼年組異常地窄，而中心化之後帶反而變寬。
本篇補上三個二階解析量。首先是對數尺度的變異數
$v(\mu;c_0)=\mathrm{Var}[\log\max(D,c_0)]$，其小 $\mu$ 展開為
$\mu(\log c_0)^2$、大 $\mu$ 展開為 $1/\mu$，兩端皆趨於零而於
$\mu=2.0365$ 取極大 $0.4805$。其次是偏誤函數的導數，由
$\frac{d}{d\mu}\Ex[g(D)]=\Ex[g(D{+}1)]-\Ex[g(D)]$ 得到精確式
$b'(\mu;c_0)=\Ex[\log(D{+}1)]-\Ex[\log\max(D,c_0)]-1/\mu$，
其彈性 $\mu b'(\mu)$ 的兩端極限為 $-1$ 與 $0$。
再者是校正的變異數膨脹因子 $(1-c_x)^2$，其中
$c_x=T^{-1}\sum_t\mu_{x,t}b'(\mu_{x,t})$。
臺灣女性 2001--2024 年資料的模擬顯示，$v(\mu;c_0)$ 對標準 Lee-Carter 逐年齡
$\alpha$ 抽樣標準差的預測，其與實測的比值在 22 個年齡上落在 $0.86$ 至 $1.15$，
三個人口規模的中位數分別為 $1.001$、$0.991$、$0.994$；
膨脹因子的預測在 $\mu>10$ 的年齡與實測吻合至小數第三位，
在 $2\le\mu\le10$ 的誤差不超過 $0.07$，而在 $\mu<0.5$ 的年齡高估約 $0.75$ 至
$1.00$，亦即一階展開的有效範圍由 $\mu$ 本身決定。
在 $\alpha_x$ 與 $\beta_x$ 的估計組成成分上，本篇的結果與 1005 報告所述的
行為相符，並進一步指出校正所付的變異數代價集中於零格主導的年齡，
而在期望死亡數大的年齡校正反而使變異數下降，其極小值為 $0.7640$
（$\mu=4.4288$）。末節寫出漸近框架所須證明的兩個命題，
並說明其技術困難在於 $\hat\kappa_t$ 的個數隨 $T$ 增長。
\par\vspace{6pt}
\textbf{關鍵詞}：Lee-Carter 模型、小人口、對數轉換、變異數閉式、
漸近框架、增長參數維度
\end{minipage}
\end{center}

\vspace{10pt}

\section{本篇的定位}\label{sec:intro}

\subsection{1005 進度報告所確立者}

1005 進度報告把標準 Lee-Carter 辨識為取平方損失的 $M$-估計量，
並依 Fisher 一致性的要求計算其估計方程在真值處的期望。
該期望不為零，且其量有閉式。記 $D_{x,t}$ 為 $x$ 歲組於 $t$ 年的死亡數、
$E_{x,t}$ 為其曝露數、$\mu_{x,t}=E_{x,t}m_{x,t}$ 為期望死亡數，
而標準 Lee-Carter 的反應變數為
$\ell_{x,t}=\log\bigl(\max(D_{x,t},c_0)/E_{x,t}\bigr)$，則
\begin{equation}\label{eq:b}
b(\mu;c_0)\;=\;\Ex\bigl[\log\max(D,c_0)\bigr]-\log\mu
\;=\;e^{-\mu}\log c_0+\sum_{d\ge1}\frac{e^{-\mu}\mu^{d}}{d!}\log d-\log\mu .
\end{equation}
由於 $\hat\alpha_x$ 依定義為 $\ell_{x,\cdot}$ 的列平均、在奇異值分解之前即已
定出，$\hat\alpha_x$ 的偏誤恰為該年齡各格 $b$ 的算術平均
$\bar b_x=T^{-1}\sum_t b(\mu_{x,t};c_0)$，此為恆等式而非漸近近似。
臺灣女性資料的驗證給出 $1.000$、$1.000$、$0.988$ 的逐年齡相關係數。

由該閉式導出的修正為：以 $\hat\mu_{x,t}$ 代入 $b$，再由 $\hat\alpha_x$ 扣除
$\bar b_x$。1005 報告並證實，校正若逐格施加會把 $\hat\mu$ 的噪音注入每一格而
使 $\hat\beta_x$ 變差；改為僅校正 $\hat\alpha_x$ 則完全不動中心化矩陣，
故與資訊加權的改善可以疊加。人數五萬時該作法使 $\alpha$ 最大偏誤由
$0.866$ 降至 $0.044$，而零歲平均餘命的偏誤由 $-1.480$ 歲降至 $+0.007$ 歲。

\subsection{所留下的二階問題}

式~\eqref{eq:b} 刻畫的是估計方程期望的偏離，屬一階量。
1005 報告的實證另留下兩項未被解釋的現象，而兩者都屬二階。

其一是標準 Lee-Carter 的 $\alpha$ 帶在幼年組異常地窄。
該報告第肆節之二記載，人數五萬時帶寬中位為 $0.218$，而 5 至 9 歲組的偏誤為
$0.866$，偏誤是帶寬的四倍；亦即該年齡幾乎每一次模擬都落在真值的同一側。
該報告把這一點歸因於「帶寬反映估計方程的變異數、中線位置反映其期望」，
但未給出該變異數的量。

其二是中心化之後帶反而變寬。同一節記載帶寬由 $0.218$ 變為 $0.277$，
且在幼年組變寬得更多。該報告把這一點歸因於「校正須代入 $\hat\mu$，
而 $b(\hat\mu)$ 不是確定性的函數」，同樣未給出該膨脹的量。

概括而論，1005 報告對偏誤給出了閉式，對變異數則只給出了機制的描述。
本篇補上後者。

\subsection{本篇的範圍與後續的分工}

本篇只處理二階解析量與漸近框架，亦即第\ref{sec:second}節與
第\ref{sec:asym}節。資料前處理的檢查程序、模型檢驗的統計量，
以及 $\hat\kappa_t$ 偏誤的校正，另於本系列的第二篇處理；
第\ref{sec:next}節列出其範圍，以便對照。

本篇所用的資料、真值、人口規模與亂數種子，均與 1005 報告完全相同
（Human Mortality Database 的臺灣女性 2001--2024 年資料，$A=22$ 個五齡組、
$T=24$ 年，真值取全國資料的 Poisson 最大概似配適，
$N\in\{10^4,5\times10^4,2\times10^5\}$，重複 $100$ 次，
種子 $20260926+1000\,i+r$），故本篇的各表可與該報告的各表逐格對照。
計算以 R 4.2.1 完成，程式為 \texttt{R/core/analytic\_variance.R} 與
\texttt{R/studies/run\_analytic\_variance.R}，不依賴外部套件。

\section{承接的文獻}\label{sec:lit}

本篇所借用的性質出自三處，此處只說明其與本篇的接點，
完整的脈絡見 1005 報告第貳節。

就\textbf{偏誤修正的施加方式}而言，本篇承接 \citet{kosmidis2010}。
該文給出一個通用的疊代演算法，只要某個模型的最大概似估計量的首階偏誤已經
導出，即可據以計算其降偏誤估計，並明言該演算法可視為一連串的疊代偏誤修正。
本篇所處理的「代入 $\hat\mu$ 所付的變異數代價」，在該文的框架中即為
疊代的不動點與其一階展開，故本篇第\ref{sec:second}節之四的推導有先例可循。
\citet{kennepagui2017} 另指出期望不偏不蘊含中位數不偏，
且前者不具再參數化不變性；該項區別屬本系列第二篇的範圍。

就\textbf{偏誤不隨樣本數消失的情形}而言，本篇承接 \citet{breslow1995}。
該文處理帶單一變異成分的廣義線性混合模型，發現一階估計量的偏誤量由變異
成分的大小決定而非由樣本數決定，並給出易算的修正因子；
\citet{linbreslow1996} 把結果推廣到多個變異成分。
本篇第\ref{sec:asym}節所要寫下的漸近敘述，其形式與該文相同：
在某一個區制下偏誤不消失，而修正在同一區制下仍然有效。

就\textbf{奇異向量的逐元誤差}而言，$\hat\beta_x$ 與 $\hat\kappa_t$ 為
中心化矩陣的奇異向量，其誤差須以逐元的矩陣擾動理論處理
（\citealp{abbe2020}），而非以子空間的角度界限處理
（\citealp{caizhang2018}）。本問題的變異數逐格不同
（第\ref{sec:second}節之二將指出 $\mathrm{Var}(\ell_{x,t})=v(\mu_{x,t})$），
恰為 \citet{zhang2022} 所處理的異質變異設定。
此一部分本篇只寫出其在漸近框架中的位置，推導留待後續。

最後，第\ref{sec:asym}節之三所述「參數維度隨 $T$ 增長」的困難，
其最接近的工具為 \citet{lahiri2023} 的巢式誤差迴歸，
該文容許區域別的迴歸係數與區域別的變異成分同時存在。

\section{二階解析量}\label{sec:second}

\subsection{記號與假設}

沿用 1005 報告的設定。令 $D\sim\text{Poisson}(\mu)$，$0<c_0\le1$ 為零格的
替代值（本文一律取 $c_0=0.5$），並記 $g_{c_0}(d)=\log\max(d,c_0)$。
由式~\eqref{eq:b}，$b(\mu;c_0)=\Ex[g_{c_0}(D)]-\log\mu$。
以下三個小節分別給出 $g_{c_0}(D)$ 的變異數、$b$ 的導數，
以及代入 $\hat\mu$ 所造成的變異數膨脹。

\subsection{對數尺度的變異數}

\begin{prop}\label{prop:v}
令 $v(\mu;c_0)=\mathrm{Var}[g_{c_0}(D)]$。則
\begin{equation}\label{eq:v}
v(\mu;c_0)=e^{-\mu}(\log c_0)^2+\sum_{d\ge1}\frac{e^{-\mu}\mu^{d}}{d!}(\log d)^2
-\Bigl[e^{-\mu}\log c_0+\sum_{d\ge1}\frac{e^{-\mu}\mu^{d}}{d!}\log d\Bigr]^{2},
\end{equation}
且其兩端的行為為
\[
v(\mu;c_0)=\mu(\log c_0)^2+O(\mu^2)\quad(\mu\to0),\qquad
v(\mu;c_0)=\frac{1}{\mu}+O(\mu^{-2})\quad(\mu\to\infty).
\]
\end{prop}

\begin{proof}
式~\eqref{eq:v} 為定義的直接展開。$\mu\to0$ 時
$P(D=0)=1-\mu+O(\mu^2)$、$P(D=1)=\mu+O(\mu^2)$，而 $g_{c_0}(0)=\log c_0$、
$g_{c_0}(1)=0$，故 $g_{c_0}(D)$ 以機率 $1-\mu+O(\mu^2)$ 取 $\log c_0$、
以機率 $\mu+O(\mu^2)$ 取 $0$。兩點分布的變異數為
$(1-\mu)\mu(\log c_0)^2+O(\mu^2)=\mu(\log c_0)^2+O(\mu^2)$。
$\mu\to\infty$ 時 $P(D=0)\to0$ 使 $\max$ 失效，
而 $\log D=\log\mu+(D-\mu)/\mu+O_p(\mu^{-1})$，
故 $\mathrm{Var}[\log D]=\mathrm{Var}(D)/\mu^2+O(\mu^{-2})=1/\mu+O(\mu^{-2})$。
\end{proof}

由 Proposition~\ref{prop:v} 可以看到此一變異數的三項特色。
首先，它在\textbf{兩端都趨於零}，故必在中間取極大；
數值求解給出極大點 $\mu=2.0365$，其值為 $0.4805$。
亦即對數尺度上最不穩定的年齡，是每年約兩人死亡的年齡，
而不是零格主導的年齡。
再來，$\mu\to0$ 的極限給出 1005 報告那個觀察的解析解釋：
零格主導的年齡其 $\ell_{x,t}$ 幾乎必然等於常數 $\log(c_0/E_{x,t})$，
重複之間幾無變動，故帶窄；但該常數與真值相差 $b(\mu)\approx\log(c_0/\mu)$，
故偏誤大。\textbf{帶窄與準確在此完全脫鉤，而脫鉤的程度可以算出來}。
其次，$\mu\to\infty$ 的極限 $1/\mu$ 恰為 Poisson 計數取對數後的標準結果，
說明式~\eqref{eq:v} 在大樣本端退化為既有的常態近似。
簡言之，這三項特色使 $v(\mu;c_0)$ 可以同時描述零格主導與計數充足兩個區域，
而這正是小人口資料跨年齡所涵蓋的範圍。

由此立得 $\hat\alpha_x$ 的變異數。

\begin{cor}\label{cor:var}
若 $\beta_x\kappa_t$ 已知，則 $\hat\alpha_x=T^{-1}\sum_t(\ell_{x,t}-\beta_x\kappa_t)$，
而各格互相獨立，故
\begin{equation}\label{eq:valpha}
\mathrm{Var}(\hat\alpha_x)=\frac{1}{T^{2}}\sum_{t}v(\mu_{x,t};c_0).
\end{equation}
\end{cor}

Corollary~\ref{cor:var} 與 1005 報告的 Corollary（$\hat\alpha_x$ 的偏誤為
$\bar b_x$）互為一對：\textbf{標準 Lee-Carter 的 $\alpha_x$，其偏誤與變異數
都是單一個期望死亡數的函數}。兩者合起來即給出均方誤差的閉式，
亦使解析信賴區間成為可能，而 1005 報告的 $\alpha_x$ 並無區間理論。

\subsection{偏誤函數的導數}

校正須代入 $\hat\mu$，故須知道 $b$ 對 $\mu$ 的敏感度。
以下的引理使該導數不必以數值微分求得。

\begin{lem}\label{lem:diff}
設 $g$ 使 $\Ex[|g(D)|]<\infty$ 且逐項微分成立。則
$\dfrac{d}{d\mu}\Ex[g(D)]=\Ex[g(D{+}1)]-\Ex[g(D)]$。
\end{lem}

\begin{proof}
$\Ex[g(D)]=\sum_{d\ge0}g(d)e^{-\mu}\mu^{d}/d!$，逐項對 $\mu$ 微分得
\[
\sum_{d\ge0}g(d)\Bigl(-\frac{e^{-\mu}\mu^{d}}{d!}
+\frac{e^{-\mu}d\,\mu^{d-1}}{d!}\Bigr)
=-\Ex[g(D)]+\sum_{d\ge1}g(d)\frac{e^{-\mu}\mu^{d-1}}{(d-1)!}
=-\Ex[g(D)]+\Ex[g(D{+}1)].
\]
\end{proof}

\begin{prop}\label{prop:bp}
設 $0<c_0\le1$。則
\begin{equation}\label{eq:bp}
b'(\mu;c_0)=\Ex\bigl[\log(D{+}1)\bigr]-\Ex\bigl[\log\max(D,c_0)\bigr]-\frac{1}{\mu},
\end{equation}
且其彈性 $\mu b'(\mu;c_0)$ 滿足
\[
\mu b'(\mu;c_0)=-1-\mu\log c_0+O(\mu^2)\quad(\mu\to0),\qquad
\mu b'(\mu;c_0)=\frac{1}{2\mu}+\frac{5}{6\mu^{2}}+O(\mu^{-3})\quad(\mu\to\infty).
\]
\end{prop}

\begin{proof}
取 $g=g_{c_0}$ 代入 Lemma~\ref{lem:diff}。因 $c_0\le1$ 而 $D+1\ge1$，
故 $g_{c_0}(D{+}1)=\log(D{+}1)$，再減去 $\log\mu$ 的導數 $1/\mu$ 即得
式~\eqref{eq:bp}。$\mu\to0$ 時 $\Ex[\log(D{+}1)]=O(\mu)$ 而
$\Ex[g_{c_0}(D)]=\log c_0+O(\mu)$，故 $\mu b'=-\mu\log c_0-1+O(\mu^2)$。
$\mu\to\infty$ 時由 1005 報告的大 $\mu$ 展開
$b(\mu)=-\frac{1}{2\mu}-\frac{5}{12\mu^{2}}+O(\mu^{-3})$ 逐項微分即得。
\end{proof}

由 Proposition~\ref{prop:bp} 可以看到兩項性質。
其一，$\mu b'(\mu)$ 的兩端極限分別為 $-1$ 與 $0$，而該量是有界的，
其值域為 $[-1,\,0.1259]$。
其二，$b'$ 有唯一的零點，位於 $\mu=2.2886$，該處 $b$ 取極小 $-0.19104$；
故 $\mu<2.2886$ 時 $b'<0$、$\mu>2.2886$ 時 $b'>0$。
此一變號點與 1005 報告的偏誤變號點 $\mu^{\ast}=0.9234$ 不同，
兩者分別是 $b$ 的零點與 $b'$ 的零點，不可混淆。
簡言之，式~\eqref{eq:bp} 把一個原本只能以數值微分近似的量化為精確式，
而其兩端的極限恰好決定了下一小節膨脹因子的上下界。

\subsection{校正的變異數膨脹因子}

修正後的估計量 $\hat\alpha^{\ast}_x$ 為不動點
$\hat\alpha^{\ast}_x=\hat\alpha_x-T^{-1}\sum_t b(\hat\mu_{x,t})$，
其中 $\hat\mu_{x,t}=E_{x,t}\exp(\hat\alpha^{\ast}_x+\hat\beta_x\hat\kappa_t)$。
以下的展開把 $\hat\beta_x\hat\kappa_t$ 視為已知，只保留 $\hat\alpha_x$ 的變動。

\begin{prop}\label{prop:infl}
記 $A=\hat\alpha_x-\alpha_x$ 與
\begin{equation}\label{eq:cx}
c_x=\frac{1}{T}\sum_{t}\mu_{x,t}\,b'(\mu_{x,t};c_0).
\end{equation}
則在 $\hat\beta_x\hat\kappa_t$ 已知的前提下，
$\hat\alpha^{\ast}_x-\alpha_x=(1-c_x)(A-\bar b_x)+o_p(A)$，
故
\begin{equation}\label{eq:infl}
\mathrm{Var}(\hat\alpha^{\ast}_x)=(1-c_x)^{2}\,\mathrm{Var}(\hat\alpha_x)+o(\cdot).
\end{equation}
\end{prop}

\begin{proof}
因 $\partial\mu_{x,t}/\partial\alpha_x=\mu_{x,t}$，一階展開給出
$T^{-1}\sum_t b(\hat\mu_{x,t})=\bar b_x+c_x A+o_p(A)$，
代入不動點的定義並整理即得。
\end{proof}

由 Proposition~\ref{prop:infl} 可以看到此一膨脹因子的三項性質，
而三者都由 Proposition~\ref{prop:bp} 的極限直接讀出。
首先，零格主導的年齡 $c_x\to-1$，膨脹因子趨於 $4$；
亦即\textbf{校正所付的變異數代價，恰好集中在它收益最大的年齡}。
再來，期望死亡數大的年齡 $c_x\to0$，膨脹因子趨於 $1$，校正不付代價。
其次，$\mu>2.2886$ 時 $b'>0$ 使 $c_x>0$，膨脹因子\textbf{小於 $1$}，
亦即校正在該區域反而使變異數下降；其極小值為 $0.7640$，發生於 $\mu=4.4288$。
最後一項與直覺相反，值得單獨指出：既有文獻一般把偏誤修正視為「以變異數換
偏誤」的交易（\citealp{firth1993}；\citealp{huber2009}），
而此處在 $\mu>2.2886$ 的年齡兩者同時改善。
其機制是校正量 $b(\hat\mu)$ 與 $\hat\alpha_x$ 正相關，故扣除它等於作了一次
部分的收縮，而該收縮在 $b$ 遞增的區域方向正確。

簡言之，式~\eqref{eq:infl} 使「校正是否划算」成為一個可以逐年齡計算的量：
把式~\eqref{eq:cx} 與 Corollary~\ref{cor:var} 合併，即得修正前後的均方誤差
之比，而該比值只依賴期望死亡數。

\subsection{三項性質的綜述}

把上述三個小節收束起來，標準 Lee-Carter 在對數尺度上的行為可以由
$\mu$ 的三個函數完全描述：偏誤 $b(\mu;c_0)$、變異數 $v(\mu;c_0)$、
以及校正的代價 $(1-c_x)^2$。三者的形狀彼此不同，其交會處即為方法的適用範圍。
圖~\ref{fig:v} 把後兩者與 $\mu b'(\mu)$ 並列。

\begin{figure}[H]
\centering
\includegraphics[width=\textwidth]{figV1_logvar.png}
\caption{三個二階解析量對期望死亡數的變化（$c_0=0.5$，橫軸取對數）。
(a) 對數尺度的變異數 $v(\mu;c_0)$，兩端趨於零而於 $\mu=2.03$ 取極大 $0.481$。
(b) 偏誤導數的彈性 $\mu b'(\mu)$，兩端極限為 $-1$ 與 $0$。
(c) 校正的變異數膨脹因子 $(1-c_x)^2$（縱軸亦取對數），
由 $4$ 遞減至極小 $0.764$（$\mu=4.43$）後回升至 $1$。}
\label{fig:v}
\end{figure}

圖~\ref{fig:v} 的設計是把三個量放在同一條 $\mu$ 軸上，以便判讀它們的相對位置。
明顯之處有二。一方面，(a) 的極大點 $\mu=2.03$ 與 (c) 的極小點 $\mu=4.43$
相當接近，亦即\textbf{變異數最大的年齡，恰好也是校正最能降低變異數的年齡}；
這與第\ref{sec:second}節之四所述的收縮機制吻合。
另一方面，(c) 在 $\mu<1$ 的區域迅速爬升至 $4$，
而 1005 報告第伍節之三所記載「資訊加權在幼年組使帶變寬得更多」正落在此處。
不明顯之處在 (b) 的右側：$\mu b'(\mu)$ 於 $\mu=4.43$ 之後緩慢回落至零，
使 (c) 在 $\mu>10$ 的區域幾乎貼齊 $1$；
該區域的校正因此既無收益亦無代價，而 1005 報告所觀察到的老年組改善有限，
其原因即在於此，而非方法在該處失效。
綜上所述，圖~\ref{fig:v} 所表現的是：三個解析量把年齡軸切成三段，
$\mu<1$ 為偏誤主導而校正昂貴，$1\le\mu\le10$ 為兩者兼具而校正划算，
$\mu>10$ 則兩者皆小而校正近乎中性。

\section{以臺灣資料檢驗}\label{sec:check}

\subsection{檢驗的設計}

第\ref{sec:second}節的三個命題，其中 Proposition~\ref{prop:v} 與
Proposition~\ref{prop:bp} 為恆等式，嚴格說來不需要檢驗；
Proposition~\ref{prop:infl} 則含一階展開，其有效範圍須以數值界定。
本節因此分三步。首先以模擬核對兩個恆等式的實作，
蓋因兩式都含無窮級數，截斷的位置可能影響結果。
再來確認 Corollary~\ref{cor:var} 的預測與實際估計流程所產生的量為同一個量，
理由與 1005 報告第肆節相同：實際的配適含奇異值分解與識別條件的正規化，
而推導把 $\beta_x\kappa_t$ 當成已知。
最後界定 Proposition~\ref{prop:infl} 的有效範圍。
資料、真值、人口規模與亂數種子均依第\ref{sec:intro}節之三所述。

\subsection{三個解析量的數值核對}

表~\ref{tab:v1} 列出三個解析量在代表性 $\mu$ 上的值。
$v$ 的核對以 $2\times10^{6}$ 次 Poisson 抽樣的樣本變異數為對照，
$b'$ 的核對以步長 $10^{-5}\mu$ 的中央差分為對照。

\begin{table}[H]
\centering
\caption{三個二階解析量的值與數值核對（$c_0=0.5$）}
\small
\begin{tabular}{rrrrrrrr}
\toprule
$\mu$ & $b(\mu)$ & $v(\mu)$ & 模擬 $v$ & 比值 & $b'(\mu)$ & $\mu b'(\mu)$ & $(1-c)^2$\\
\midrule
0.02   & $3.2327$  & 0.0096 & 0.0095 & 1.006 & $-49.3069$ & $-0.9861$ & 3.945\\
0.05   & $2.3372$  & 0.0240 & 0.0241 & 0.998 & $-19.3072$ & $-0.9654$ & 3.863\\
0.10   & $1.6787$  & 0.0479 & 0.0479 & 0.999 & $-9.3082$  & $-0.9308$ & 3.728\\
0.30   & $0.7176$  & 0.1400 & 0.1401 & 0.999 & $-2.6513$  & $-0.7954$ & 3.223\\
0.50   & $0.3416$  & 0.2230 & 0.2230 & 1.000 & $-1.3346$  & $-0.6673$ & 2.780\\
0.6931 & $0.1416$  & 0.2921 & 0.2925 & 0.999 & $-0.7981$  & $-0.5532$ & 2.412\\
0.9234 & $0.0000$  & 0.3589 & 0.3585 & 1.001 & $-0.4667$  & $-0.4309$ & 2.048\\
1.00   & $-0.0329$ & 0.3774 & 0.3772 & 1.000 & $-0.3937$  & $-0.3937$ & 1.943\\
2.00   & $-0.1874$ & \textbf{0.4805} & 0.4807 & 1.000 & $-0.0278$ & $-0.0555$ & 1.114\\
5.00   & $-0.1199$ & 0.2868 & 0.2866 & 1.001 & $0.0246$ & $0.1228$ & \textbf{0.770}\\
10.0   & $-0.0552$ & 0.1206 & 0.1205 & 1.001 & $0.0062$ & $0.0618$ & 0.880\\
20.0   & $-0.0262$ & 0.0543 & 0.0543 & 1.000 & $0.0014$ & $0.0274$ & 0.946\\
50.0   & $-0.0102$ & 0.0206 & 0.0207 & 0.999 & $0.0002$ & $0.0104$ & 0.979\\
100    & $-0.0050$ & 0.0102 & 0.0102 & 1.000 & $0.0001$ & $0.0051$ & 0.990\\
300    & $-0.0017$ & 0.0034 & 0.0033 & 1.001 & $0.0000$ & $0.0017$ & 0.997\\
\bottomrule
\end{tabular}

\vspace{3pt}
{\footnotesize $b'$ 與中央差分的比值在全部 15 個 $\mu$ 上皆為 $1.0000$，
故不另列。$\mu=0.6931$ 為中位數失效門檻 $\ln2$，
$\mu=0.9234$ 為偏誤變號點 $\mu^{\ast}$。粗體為該欄的極值。}
\label{tab:v1}
\end{table}

表~\ref{tab:v1} 的設計是在同一組 $\mu$ 上並列三個解析量與其數值對照，
以確認實作的截斷不影響結果。
從表中可以看到，$v$ 的比值在 15 個 $\mu$ 上落在 $0.998$ 至 $1.006$，
而 $b'$ 與中央差分的比值一律為 $1.0000$，兩者都在蒙地卡羅誤差與浮點誤差的
量級之內，故兩式的實作正確。
特別明顯的是兩端的展開：$\mu=0.02$ 時 $v=0.0096$ 而
$\mu(\log c_0)^2=0.0096$，$\mu=300$ 時 $v=0.0034$ 而 $1/\mu=0.0033$，
與 Proposition~\ref{prop:v} 的兩個極限吻合；
$\mu=0.02$ 時 $\mu b'=-0.9861$ 而 $-1-\mu\log c_0=-0.9861$，
$\mu=300$ 時 $\mu b'=0.001676$ 而 $\frac{1}{2\mu}+\frac{5}{6\mu^2}=0.001676$，
與 Proposition~\ref{prop:bp} 的兩個展開吻合至小數第六位。
不明顯之處在中段：$\mu$ 由 $2$ 至 $5$ 之間，$v$ 由極大 $0.4805$ 下降至
$0.2868$ 而 $(1-c)^2$ 由 $1.114$ 下降至極小 $0.770$，兩者的轉折並不在同一點，
這是因為前者由 $b$ 的二階矩決定而後者由 $b$ 的一階導數決定，
兩者本無須同時轉折。
綜上所述，表~\ref{tab:v1} 所表現的是三個解析量的實作與其理論極限一致，
可作為後續逐年齡比較的基礎。

\subsection{逐年齡的變異數預測}

表~\ref{tab:v2} 以 $N=5\times10^{4}$ 為例，列出 Corollary~\ref{cor:var} 的
預測與 100 次重複的實測抽樣標準差。

\begin{table}[H]
\centering
\caption{$\hat\alpha_x$ 抽樣標準差的預測與實測（$N=5\times10^{4}$，重複 100 次）}
\small
\begin{tabular}{lrrrrrrr}
\toprule
年齡組 & $\mu$ 中位 & 閉式 sd & 實測 sd & 比值 & $c_x$ & 閉式膨脹 & 實測膨脹\\
\midrule
0        & 1.111  & 0.1299 & 0.1407 & 1.083 & $-0.329$ & 1.766 & 1.835\\
1--4     & 0.283  & 0.0778 & 0.0740 & 0.951 & $-0.786$ & 3.191 & 2.477\\
5--9     & 0.241  & 0.0707 & 0.0752 & 1.064 & $-0.825$ & 3.331 & 2.543\\
10--14   & 0.254  & 0.0716 & 0.0713 & 0.996 & $-0.820$ & 3.314 & 2.527\\
15--19   & 0.497  & 0.0978 & 0.0893 & 0.914 & $-0.656$ & 2.742 & 2.254\\
20--24   & 0.829  & 0.1189 & 0.1019 & 0.857 & $-0.467$ & 2.151 & 2.040\\
25--29   & 1.204  & 0.1324 & 0.1518 & 1.146 & $-0.290$ & 1.663 & 1.654\\
30--34   & 1.720  & 0.1400 & 0.1527 & 1.091 & $-0.116$ & 1.245 & 1.564\\
35--39   & 2.554  & 0.1383 & 0.1361 & 0.984 & $0.042$  & 0.918 & 1.312\\
40--44   & 4.847  & 0.1102 & 0.1097 & 0.995 & $0.122$  & 0.771 & 0.851\\
45--49   & 7.194  & 0.0868 & 0.0856 & 0.986 & $0.090$  & 0.827 & 0.807\\
50--54   & 9.921  & 0.0708 & 0.0628 & 0.888 & $0.062$  & 0.881 & 0.871\\
55--59   & 14.641 & 0.0562 & 0.0513 & 0.914 & $0.039$  & 0.925 & 0.921\\
60--64   & 22.595 & 0.0442 & 0.0454 & 1.028 & $0.024$  & 0.953 & 0.950\\
65--69   & 32.653 & 0.0363 & 0.0324 & 0.893 & $0.016$  & 0.968 & 0.969\\
70--74   & 46.185 & 0.0303 & 0.0311 & 1.026 & $0.011$  & 0.978 & 0.978\\
75--79   & 46.091 & 0.0303 & 0.0289 & 0.952 & $0.011$  & 0.978 & 0.978\\
80--84   & 64.868 & 0.0255 & 0.0231 & 0.906 & $0.008$  & 0.985 & 0.984\\
85--89   & 72.332 & 0.0241 & 0.0247 & 1.023 & $0.007$  & 0.986 & 0.986\\
90--94   & 53.280 & 0.0282 & 0.0267 & 0.947 & $0.010$  & 0.981 & 0.981\\
95--99   & 22.967 & 0.0440 & 0.0452 & 1.027 & $0.023$  & 0.954 & 0.952\\
$100^{+}$& 5.035  & 0.1086 & 0.1097 & 1.010 & $0.122$  & 0.771 & 0.797\\
\midrule
中位數    &        & 0.0712 & 0.0727 & 0.991 &          & 0.979 & 0.983\\
\bottomrule
\end{tabular}
\label{tab:v2}
\end{table}

表~\ref{tab:v2} 的設計是在固定人口規模下逐年齡並列兩組量：
Corollary~\ref{cor:var} 對標準 Lee-Carter 抽樣標準差的預測與其實測值，
以及 Proposition~\ref{prop:infl} 對校正後變異數膨脹的預測與其實測值。
從表中可以看到，標準差的比值在 22 個年齡上落在 $0.857$ 至 $1.146$，
中位數為 $0.991$，而其離散與 100 次重複所能給出的精度相當；
亦即\textbf{標準 Lee-Carter 的 $\alpha$ 變異數亦可由單一個期望死亡數的函數
預測}，與 1005 報告對偏誤所得的結論同型。
特別明顯的是 $\mu$ 由 $0.24$ 升至 $72$ 時，閉式 sd 由 $0.0707$ 先升至
$0.1400$（$\mu=1.72$）再降至 $0.0241$，實測完全跟上這個非單調的型態，
而該型態正是 Proposition~\ref{prop:v} 所預言者。
膨脹因子的部分，$\mu>10$ 的九個年齡（55--59 歲至 95--99 歲）
預測與實測吻合至小數第三位，例如 80--84 歲為 $0.985$ 對 $0.984$、
85--89 歲為 $0.986$ 對 $0.986$。

不明顯之處集中在 $\mu<1$ 的年齡，而該處閉式\textbf{高估}膨脹：
5--9 歲組預測 $3.331$ 而實測 $2.543$，1--4 歲組預測 $3.191$ 而實測 $2.477$。
此一落差來自 Proposition~\ref{prop:infl} 的一階展開在該處失效。
就本文所能判斷，原因有二。一方面 $b$ 在 $\mu$ 小時高度凸
（由 Proposition~\ref{prop:bp}，$b'\approx-1/\mu$ 而 $b''\approx1/\mu^2$），
故一階項不足以描述 $b(\hat\mu)$ 的變動。
另一方面 $\hat\mu$ 在該處受零格替代的牽制，其變動並不如展開式所假設的那樣
自由。這構成本篇推導的一項限制，而該限制的位置明確：
膨脹因子的預測只在 $\mu\gtrsim1$ 可用。
另有一處實測高於預測，即 30--34 歲與 35--39 歲兩組
（$1.245$ 對 $1.564$、$0.918$ 對 $1.312$），
該處 $c_x$ 由負轉正而 $v$ 接近極大，兩項變化同時發生，
故一階展開的誤差在該窄帶內換號；此一現象尚未查明，列為後續工作。

綜上所述，表~\ref{tab:v2} 所表現的是：$\alpha$ 的變異數與偏誤同樣可由閉式
預測，而校正的代價在 $\mu\gtrsim1$ 的年齡亦可由閉式預測；
至於 $\mu<1$ 的年齡，閉式只能給出代價的上界。

\subsection{膨脹因子的有效範圍}

為把上一小節的落差寫成可用的判準，表~\ref{tab:v3} 把三個人口規模的
22 個年齡依期望死亡數分段彙總。

\begin{table}[H]
\centering
\caption{膨脹因子的預測誤差依期望死亡數分段（重複 100 次）}
\small
\begin{tabular}{lrrrrr}
\toprule
$\mu$ 區間 & 年齡數 & sd 比值 & 閉式膨脹 & 實測膨脹 & 預測誤差\\
\midrule
\multicolumn{6}{l}{\textit{$N=10^{4}$}}\\
$\mu<0.5$   & 8 & 1.027 & 3.630 & 2.630 & $+0.999$\\
$0.5\le\mu<1$ & 2 & 0.970 & 2.344 & 2.119 & $+0.225$\\
$1\le\mu<2$   & 3 & 1.106 & 1.456 & 1.706 & $-0.240$\\
$2\le\mu<10$  & 6 & 1.001 & 0.832 & 0.853 & $-0.024$\\
$\mu\ge10$    & 3 & 0.988 & 0.913 & 0.905 & $+0.008$\\
\midrule
\multicolumn{6}{l}{\textit{$N=5\times10^{4}$}}\\
$\mu<0.5$   & 4 & 0.973 & 3.252 & 2.502 & $+0.750$\\
$0.5\le\mu<1$ & 1 & 0.857 & 2.151 & 2.040 & $+0.111$\\
$1\le\mu<2$   & 3 & 1.091 & 1.663 & 1.654 & $-0.069$\\
$2\le\mu<10$  & 5 & 0.986 & 0.827 & 0.851 & $-0.026$\\
$\mu\ge10$    & 9 & 0.952 & 0.978 & 0.978 & $+0.000$\\
\midrule
\multicolumn{6}{l}{\textit{$N=2\times10^{5}$}}\\
$0.5\le\mu<1$ & 1 & 0.953 & 1.932 & 1.860 & $+0.072$\\
$1\le\mu<2$   & 3 & 1.163 & 1.697 & 1.717 & $-0.019$\\
$2\le\mu<10$  & 4 & 0.978 & 0.791 & 0.839 & $-0.066$\\
$\mu\ge10$    & 14 & 1.003 & 0.989 & 0.989 & $+0.000$\\
\bottomrule
\end{tabular}

\vspace{3pt}
{\footnotesize 各欄為該段內各年齡的中位數。$N=2\times10^{5}$ 時無 $\mu<0.5$
的年齡，故該列從缺。}
\label{tab:v3}
\end{table}

表~\ref{tab:v3} 的設計是把年齡依期望死亡數分成五段，並在三個人口規模下
各自彙總，以便判斷一階展開的失效是由 $\mu$ 決定還是由 $N$ 決定。
從表中可以看到，預測誤差的型態在三個規模下完全一致，
且其大小由 $\mu$ 所屬的區間決定而非由 $N$ 決定：
$\mu\ge10$ 的段，三個規模的誤差分別為 $+0.008$、$+0.000$、$+0.000$；
$2\le\mu<10$ 的段為 $-0.024$、$-0.026$、$-0.066$；
$\mu<0.5$ 的段則為 $+0.999$ 與 $+0.750$。
此一一致性與 Proposition~\ref{prop:infl} 的性質吻合，
蓋因式~\eqref{eq:cx} 中 $N$ 只透過 $\mu_{x,t}$ 進入，並不單獨出現。
特別明顯的是 sd 比值這一欄：它在全部 14 個分段上都落在 $0.857$ 至 $1.163$，
不隨 $\mu$ 區間而系統性變化，可見 Corollary~\ref{cor:var} 的有效範圍
涵蓋全部年齡，與膨脹因子的情形不同。
不明顯之處在 $1\le\mu<2$ 這一段：三個規模的誤差為 $-0.240$、$-0.069$、
$-0.019$，隨 $N$ 增大而縮小，是全表唯一不符合「只由 $\mu$ 決定」的位置。
就本文所能判斷，該段恰在 $c_x$ 變號與 $v$ 取極大的交會處，
一階展開的餘項在此最大，而餘項本身含 $\hat\mu$ 的抽樣誤差，
故仍受 $N$ 影響；此一解釋尚未驗證，列為後續工作。

綜上所述，表~\ref{tab:v3} 所表現的是一條可用的判準：
$\mu\ge2$ 的年齡可直接採用式~\eqref{eq:infl} 的預測，
$1\le\mu<2$ 的年齡其預測有約 $0.02$ 至 $0.24$ 的誤差，
$\mu<1$ 的年齡則只能把式~\eqref{eq:infl} 當成代價的上界。
而在小人口的應用中，$\mu<1$ 的年齡恰好就是幼年組，
故此一限制與 1005 報告第伍節之三所記載「資訊加權在幼年組的效果最難預測」
指向同一處。

\section{漸近框架}\label{sec:asym}

\subsection{兩個區制}

1005 報告指出 $b$ 不隨觀測年數 $T$ 增加而消失，但未寫出對應的漸近敘述。
此處先把兩個區制分開，蓋因兩者對同一個估計量給出相反的結論。

\textbf{小區域區制。} 期望死亡數有界，即存在 $0<\underline\mu\le\bar\mu<\infty$
使 $\underline\mu\le\mu_{x,t}\le\bar\mu$ 對所有 $(x,t)$ 成立，
而 $T\to\infty$。此一區制描述的是鄉鎮市區：人口規模固定，
而觀測年數可以逐年累積。

\textbf{大母體區制。} $N\to\infty$ 使 $\mu_{x,t}\to\infty$，$T$ 固定。
此一區制描述的是全國或大縣市，亦即既有 Lee-Carter 文獻的標準設定。

兩個區制的差別在於 $b$ 的命運。由 1005 報告的大 $\mu$ 展開
$b(\mu)=-\frac{1}{2\mu}+O(\mu^{-2})$，大母體區制下 $b\to0$，
故標準 Lee-Carter 與本文的修正皆一致，兩者的差別為 $O(N^{-1})$。
小區域區制下 $b$ 不動，故 $\bar b_x$ 不隨 $T$ 消失。
這正是 1005 報告所處理的情形，也是 \citet{breslow1995} 所處理情形的同型：
該文的偏誤量由變異成分決定，本文的偏誤量由 $c_0$ 與期望死亡數決定，
兩者都與樣本數無關。

\subsection{小區域區制下的命題}

以下兩個命題是本篇所要證明而尚未完成者，此處先寫下其敘述，
以便下一階段逐項補上證明。

\textbf{待證命題一（一致性的分離）。}
在小區域區制下，若 $\beta_x\kappa_t$ 已知，則當 $T\to\infty$ 時
\[
\hat\alpha_x\;\xrightarrow{\;a.s.\;}\;\alpha_x+\lim_{T\to\infty}\bar b_x,
\qquad
\hat\alpha^{\ast}_x\;\xrightarrow{\;a.s.\;}\;\alpha_x .
\]
亦即標準 Lee-Carter 在該區制下\textbf{不一致}，而修正後的估計量一致。
此一敘述是本文方法存在的理由，其證明只需強大數法則與 $b$ 的連續性，
困難不在這裡。

\textbf{待證命題二（漸近常態與變異數）。}
在同一前提下，
\[
\sqrt{T}\bigl(\hat\alpha^{\ast}_x-\alpha_x\bigr)
\;\xrightarrow{\;d\;}\;N\bigl(0,\;\sigma^{2}_x\bigr),
\qquad
\sigma^{2}_x=(1-c_x)^{2}\,\lim_{T\to\infty}\frac{1}{T}\sum_t v(\mu_{x,t};c_0),
\]
其中 $c_x$ 如式~\eqref{eq:cx}。第\ref{sec:check}節之三已給出該式在
$\mu\gtrsim1$ 時的數值支持。此式一旦成立，即可直接寫出 $\hat\alpha^{\ast}_x$
的解析信賴區間，而 1005 報告第陸節的兩步校正區間即可作為對照。

\subsection{增長參數維度的困難}

上述兩個命題的前提是 $\beta_x\kappa_t$ 已知，而實際上兩者都是估計出來的。
把它們放回去之後，困難的性質就改變了：$\hat\kappa_t$ 共有 $T$ 個，
其個數隨 $T$ 一同增長，故此處不是標準的 $M$-估計漸近，
而是增長參數維度的問題。這一點與古典的 incidental parameters 情形同型，
而在小區域估計的文獻中，\citet{lahiri2023} 的巢式誤差迴歸正是處理
「區域別係數與區域別變異成分同時隨區域數增長」的設定，
為目前最接近的工具。

另一項困難在 $\hat\beta_x$。它與 $\hat\kappa_t$ 同為中心化矩陣的奇異向量，
故其誤差須以逐元的矩陣擾動理論處理（\citealp{abbe2020}），
而非以子空間的角度界限處理（\citealp{caizhang2018}），
蓋因後者給出的是整個子空間的界限，無法指出各年齡 $\hat\beta_x$ 各自偏離的量。
本問題另有一項特殊之處：由 Proposition~\ref{prop:v}，
中心化矩陣各格的變異數 $v(\mu_{x,t})$ 逐格不同且相差可達百倍，
故須採異質變異版本的結果（\citealp{zhang2022}）。
1005 報告第參節之五已把 $\hat\beta_x$ 的誤差分解為確定性擾動
$\Delta_{x,t}=b(\mu_{x,t})-\bar b_x$ 與隨機部分，前者佔 $9.2\%$；
逐元理論所要給出的即是後者的閉式。

概括而論，漸近框架的工作可以分成兩層：
把 $\beta_x\kappa_t$ 當已知的那一層，其證明的障礙不大，應先完成；
把兩者放回去的那一層，則須借助增長參數維度與逐元擾動兩套工具，
工作量遠大於前者。本篇只完成了前一層所需的解析量。

\section{小結與後續}\label{sec:next}

\subsection{本篇的結果}

本篇在 1005 報告的一階閉式之外，補上三個二階解析量，並以臺灣資料檢驗。

在 $\alpha_x$ 與 $\beta_x$ 的估計組成成分上，本篇的結果與 1005 報告所述的
行為相符，並把其中兩項未解釋的現象化為可計算的量。
就 $\alpha_x$ 而言，其變異數為 $T^{-2}\sum_t v(\mu_{x,t};c_0)$，
而 $v$ 在 $\mu\to0$ 時趨於零；這說明了「帶窄而偏誤大」何以能夠並存，
並使 $\alpha_x$ 的均方誤差成為閉式。
就校正的代價而言，其膨脹因子為 $(1-c_x)^2$，由 $c_x$ 的兩端極限
$-1$ 與 $0$ 得知該因子介於 $4$ 與 $1$ 之間，
且在 $\mu>2.2886$ 的年齡小於 $1$。
據此可知校正的代價集中於零格主導的年齡，
而在期望死亡數大的年齡校正同時改善偏誤與變異數。

然而，本篇同時證實該預測有明確的有效範圍。
膨脹因子的一階展開在 $\mu\ge2$ 的年齡與實測吻合至小數第二位以內，
在 $\mu<0.5$ 的年齡則高估 $0.75$ 至 $1.00$，
而後者恰為小人口應用中最關鍵的幼年組。
據此亦可確認二階展開或 $\hat\mu$ 分布的直接計算為必要的下一步，
而非可有可無的精緻化。

\subsection{下一階段的工作}

依優先序列出五項，前三項屬本系列第二篇。

首先是 $\hat\kappa_t$ 的確定性偏誤。1005 報告第伍節之七證實漂移偏誤由
$0.3339$ 降至 $0.1407$ 的改善全部來自資訊加權，僅校正 $\alpha$ 完全不動漂移。
但中心化矩陣帶有確定性擾動 $\Delta_{x,t}$，而漂移為
$(\hat\kappa_T-\hat\kappa_1)/(T-1)$，是一個差分，會濾掉隨機部分而保留確定性
部分；故即使 $\Delta$ 只佔 $\hat\beta$ 總誤差的 $9.2\%$，
它對漂移的影響可能遠大於此。若此一假設成立，
$\kappa$ 的校正式即為把 $\Delta$ 投影到 $\hat\beta$ 方向後扣除，
而該校正與 $\alpha$ 的校正正交。

其次是資料前處理的檢查程序。式~\eqref{eq:b} 只需要 $E_{x,t}$ 與一個近似的
$m_{x,t}$，兩者在配適之前都已知，故標準 Lee-Carter 將產生多少偏誤可以在
配適之前算出。另須檢查三件事：零格個數是否與 $e^{-\hat\mu}$ 相符
（若有真正的零膨脹，則混合權重已知這個前提失效）；
是否有過度離散（\citealp{djeundje2010}，若有則須改用負二項的 $b$）；
以及中心化矩陣的奇異值比是否高於 \citet{benaych2012} 的門檻
（若否則任何方法都無法回復 $\beta$）。

再者是模型檢驗的統計量。校正的目標是使估計方程在真值處的期望為零，
而其樣本版本在校正後應接近零，其虛無分布由 $v(\mu;c_0)$ 給出；
此即一個有閉式虛無分布的逐年齡檢定，可用以定出校正已奏效的年齡。
\citet{kennepagui2017} 的中位數偏誤修正另提供一個不同的標的，
對平均餘命這類非線性函數的標的量可能更合適。

此外是本篇所留下的兩項：膨脹因子的二階展開，
以及 $1\le\mu<2$ 這一段預測誤差隨 $N$ 變化的原因。
最後是把 $\beta_x\kappa_t$ 放回漸近框架，此項工作量最大，列為最後。

\section*{參考文獻}
\addcontentsline{toc}{section}{參考文獻}

\renewcommand{\bibsection}{}
\setlength{\bibsep}{2pt}
\bibliographystyle{apalike}
\bibliography{references}

\end{document}
